\section{Displays with break bits}\label{sec:displays}

This section constructs the shown sets of the label: at every time, the band groups of all heads together with a \emph{display}
of the frozen groups of suspended frames. \Cref{sec:displays-why} explains informally why the display is random.
\Cref{sec:displays-one} defines displays at one time, \cref{sec:displays-entries} the entries of a frame system and the composed
display of all heads, and \cref{sec:displays-D} proves that the display is small (Lemma~D). The last three subsections control the
case (ABOVE) of Lemma~V: the display above a group is a function of a \emph{view} (\cref{sec:displays-view}), the view changes
only at standard events (Proposition~W, \cref{sec:displays-W}), and a hop without a standard event above its group is exposed only
through a broken bit (\cref{sec:displays-exposure}).

\subsection{Why break bits}\label{sec:displays-why}

This subsection is an informal explanation; nothing below depends on it.

By Lemma~V, an upper hop $(p,u)$ of a head with deciding group~$d$ is violated only through truncation, a wrap, or the case
(ABOVE): the shown groups above~$d$ differ at $p$ and at~$u$. The shown sets must contain the band groups of all heads, since the
deciding group of an aligned hop is the band group of its head at both ends. Band groups alone are not enough. When a head enters
a frame (a \emph{push}), the band of the suspended segment freezes, and when the frame ends (a \emph{pop}) the segment resumes at
its frozen band. If only the current band groups were shown, a pop at a high band would add a group to the shown set, and hops of
other heads spanning the pop, with lower deciding groups, would fall into (ABOVE). The display therefore also shows frozen groups,
the \emph{tops} of the \emph{entries} of \cref{def:entries}.

The question is which tops to show. To keep the shown groups above~$d$ independent of what happens below~$d$, an entry should be
displayed as soon as its \emph{window}, the range of groups between the current activity inside it and its top, contains a shown
value. With this rule the display above~$d$ is a function of the configuration above~$d$ (\cref{lem:view-static}), and that part of
the configuration changes only at standard events, which are paid for (\cref{prop:W}). But the rule cascades: an entry displayed
because its window contains a displayed top makes the next entry whose window contains its own top displayed, and so on up the
stacks of suspended frames. When the stacks are deep, so is the cascade, and no bounded number of digits can carry the display.

Random \emph{break bits} cut the cascade. Every entry carries a bit that is~$0$, making the entry \emph{broken}, with probability
$1/Q$, independently of the other bits, and the cascade passes only through unbroken entries. From a value~$y$ it moves to the
tops of the \emph{candidates} at~$y$, at most one per head, and it stops as soon as all the candidates it meets are broken; each
step reveals at most $n$ new bits, so long cascades are unlikely (Lemma~D, \cref{lem:D}). The display above~$d$ still depends only
on the configuration above~$d$, except through the candidates at~$d$: when these are unbroken and $d$ is shown, they are displayed
exactly as under the cascading rule. So a hop can be exposed only through one of at most $2n$ candidates, $n$ at each end, being
broken, which has probability at most $2n/Q$ (\cref{prop:exposure}).

\subsection{Displays at one time}\label{sec:displays-one}

A \emph{configuration} consists of a finite set~$I$ of \emph{entries}, each with a \emph{head} $h(e)\in\{0,\dots,n-1\}$, a
\emph{bottom} $m(e)\in\N$ and a \emph{top} $c(e)\in\N$; a set $U\subseteq I$ of \emph{unbroken} entries; and a finite set
$\Sigma\subseteq\N$ of \emph{seeds}. The \emph{window} of~$e$ is the set of groups $[m(e),c(e)]$, which is empty if $m(e)>c(e)$. An
entry~$e$ is a \emph{candidate} at a value~$y$ if $m(e)\le y<c(e)$.

\begin{definition}[Displays]\label{def:display}
\formal{Display.lean}{Disp, OneCandidate, dispVals, shownSet}
The \emph{display} $\Disp$ of a configuration is the least set of entries with the following three properties.
\begin{itemize}[leftmargin=3em]
\item[\textup{(Z)}] If $m(e)=0$, then $e\in\Disp$.
\item[\textup{(S)}] If $e\in U$, $z\in\Sigma$ and $m(e)\le z\le c(e)$, then $e\in\Disp$.
\item[\textup{(C)}] If $e\in U$, $e'\in\Disp$ and $m(e)\le c(e')\le c(e)$, then $e\in\Disp$.
\end{itemize}
The \emph{displayed values} are the seeds and the tops of the displayed entries (\lean{dispVals}), and the \emph{shown set} is
$\Ok=\{0\}\cup\Sigma\cup\{c(e):e\in\Disp\}$ (\lean{shownSet}). The configuration has the \emph{candidate property} if for every
head~$v$ and every value~$y$ at most one entry~$e$ with $h(e)=v$ is a candidate at~$y$ (\lean{OneCandidate}).
\end{definition}

Rule~(Z) displays the \emph{fresh} entries, whose window starts at~$0$, whatever their bits; rules (S) and~(C) are the cascade,
through unbroken entries only. In Lean, $\Disp$ is an inductive predicate with one constructor for each rule. An entry with an
empty window is never displayed. The only structural fact about stacks that Lemma~D uses is a consequence of the candidate
property: every value has at most $n$ candidates (\lean{cands}, \lean{card_cands_le}).

\subsection{Entries and the composed display}\label{sec:displays-entries}

Let $\mathcal S$ be a frame system on $[0,T)$ with walk~$x$ (\cref{def:framesystem}), and fix a grid $(h_0,r)$. A frame $F=(f_1,f_2)$
is the interval $[f_1,f_2)$, and $F'\subseteq F$ denotes inclusion of intervals. We write $\inner(t)$ for the innermost frame
of~$t$, and $s_F(t)$ and $M_F(t)$ for the segment start and the running maximum of a frame~$F$ at~$t$
(\lean{FrameSystem.segStart}, \lean{FrameSystem.runMax}); the band is $\beta(t)=\band(M_{\inner(t)}(t))$. The \emph{segment of~$F$
current at~$t$} is the segment of~$F$ that starts at $s_F(t)$.

\begin{definition}[Entries, tops and bottoms]\label{def:entries}
\formal{Entries.lean}{FrameSystem.parent, FrameSystem.IsEntry, FrameSystem.headGroup, FrameSystem.entryTop, FrameSystem.innerEntries, FrameSystem.entryBot}
\begin{enumerate}[label=(\alph*)]
\item The \emph{parent} $\pi(F)$ of a frame~$F$ is a shortest frame strictly containing~$F$, and the root if there is none. For a
nonempty frame other than the root, $\pi(F)$ lies inside every frame strictly containing~$F$ (\lean{FrameSystem.parent_within}).
\item A frame~$F$ is an \emph{entry} if it is nonempty, it is not the root, and its \emph{parent segment}, the segment of $\pi(F)$
current at~$f_1$, is innermost again after~$F$ ends: some time $t'\ge f_2$ is an own time of~$\pi(F)$ with
$s_{\pi(F)}(t')=s_{\pi(F)}(f_1)$. The entries containing~$t$ form the \emph{stack} at~$t$.
\item The \emph{band group} at~$t$ is $g(t)=G(\beta(t))$.
\item The \emph{top} of a frame~$F$ is $c_F=G\bigl(\band(M_{\pi(F)}(f_1))\bigr)$, the group of the band at which the parent segment is
suspended.
\item The \emph{bottom} of~$F$ at~$t$ is
\[
  m_F(t)=\max\bigl(\{g(t)\}\cup\{c_{F'}:F'\text{ an entry with }F'\subseteq F,\ F'\ne F,\ t\in F'\}\bigr),
\]
and $[m_F(t),c_F]$ is the \emph{window} of~$F$ at~$t$.
\end{enumerate}
\end{definition}

An entry is a \emph{suspend child}: a frame after which its parent segment resumes. For a nonempty frame other than the root that
is not an entry, the parent segment is never innermost again from the start of the frame on. While an entry is open, the running
maximum of its parent does not change, because the times of the entry are not own times of the parent (\cref{lem:dynamics}); so its
top is the group of the frozen band of the parent segment, the band group with which that segment resumes, unless its running
maximum rises at that time. The
bottom starts at the band group of the head and rises with the tops of the entries stacked inside; the window becomes empty when
the activity inside~$F$ rises above~$c_F$. Lean chooses $\pi(F)$ among the frames of minimal length; only nonempty frames are
entries, and their parents are determined.

\begin{definition}[The composed display]\label{def:composed}
\formal{Composition.lean}{EntryIdx, dispBot, dispTop, seeds}
Let $\mathcal S_0,\dots,\mathcal S_{n-1}$ be frame systems on $[0,T)$, one for each head, with a common grid $(h_0,r)$, and write $g_v$, $c^{(v)}_F$
and $m^{(v)}_F$ for the band groups, tops and bottoms of~$\mathcal S_v$. The \emph{composed display} at time~$t$ is the configuration whose
entries are the pairs $(v,F)$ of a head~$v$ and a frame~$F$ of~$\mathcal S_v$ (\lean{EntryIdx}); $\iota$ denotes their number. The entry
$(v,F)$ has head~$v$ and top $c^{(v)}_F$; its bottom is $m^{(v)}_F(t)$ if $F$ is an entry of~$\mathcal S_v$ containing~$t$, and $c^{(v)}_F+1$
otherwise. The seeds are the band groups $g_v(t)$ of all heads. Bits $\xi$, one value $\xi_{(v,F)}\in\{0,\dots,Q-1\}$ for every
entry, decide the unbroken entries: those with $\xi_{(v,F)}\ne0$.
\end{definition}

The bits do not depend on~$t$. A pair $(v,F)$ whose frame is not on the stack of~$v$ at~$t$ has an empty window, so it is neither a
candidate nor displayed at~$t$; the composed display lives on the stacks of the heads.

\begin{lemma}[One candidate per head]\label{lem:onecandidate}
\formal{Entries.lean}{FrameSystem.entryTop_le_entryBot, FrameSystem.entry_candidate_unique}
\formal{Composition.lean}{onStack_of_cand, card_seeds_le, oneCandidate_disp}
\begin{enumerate}[label=(\alph*)]
\item If $F'\subseteq F$ are distinct entries of~$\mathcal S$ containing~$t$, then $c_{F'}\le m_F(t)$.
\item If $F$ and~$F'$ are entries of~$\mathcal S$ containing~$t$ with $m_F(t)\le y<c_F$ and $m_{F'}(t)\le y<c_{F'}$, then $F=F'$.
\item At every time, the composed display has the candidate property and at most $n$ seeds.
\end{enumerate}
\end{lemma}

\begin{outline}
(a) holds by the definition of the bottom as a maximum (\lean{FrameSystem.entryTop_le_entryBot}). (b) Two frames containing~$t$
are nested, by laminarity (\cref{def:framesystem}). If $F'\subseteq F$ and $F'\ne F$, then $y<c_{F'}\le m_F(t)\le y$ by~(a), a
contradiction, and symmetrically (\lean{FrameSystem.entry_candidate_unique}). (c) A candidate has a nonempty window, so its frame
is an entry of its head's stack and its bottom is the true bottom (\lean{onStack_of_cand}); apply~(b) within one head
(\lean{oneCandidate_disp}). The seeds are the $n$ band groups (\lean{card_seeds_le}).
\end{outline}

Along the stack of a head, inner windows lie below outer bottoms. This is why the candidates at any value are few: at most one per
head, and at most $n$ in all.

\subsection{Lemma D: the display is small}\label{sec:displays-D}

Fix a configuration with entries~$I$ in which every value has at most $n$ candidates; this holds for every configuration with
$n$ heads and the candidate property (\lean{card_cands_le}). Let a bit vector $\xi\in\{0,\dots,Q-1\}^I$ decide the unbroken entries,
$U=\{e:\xi_e\ne0\}$. Probability bounds are counts of bit vectors: a bound ``for at most $N$ bit vectors'' means probability at
most $N/Q^{\card I}$ for uniformly random bits.

The \emph{front} from a value~$z$ is the sequence $y_0=z$ and
\[
  y_{i+1}=\min\{c(e):e\text{ an unbroken candidate at }y_i\},
\]
with $y_{i+1}=y_i$ if there is no unbroken candidate at~$y_i$ (\lean{front}, \lean{frontStep}). The front \emph{runs} at step~$i$
if some candidate at~$y_i$ is unbroken (\lean{Running}); then $y_i<y_{i+1}$ (\lean{lt_frontStep}). A front that does not run at
step~$i$ stays at~$y_i$ from then on (\lean{front_add_of_not_running}), so the steps at which it runs form an initial segment
(\lean{Running.anti}). Let $\Sigma_0$ be the set of the tops $c(e)\ge1$ of the entries with $m(e)=0$ (\lean{zeroTops}); these
entries are candidates at~$0$, so $\card{\Sigma_0}\le n$ (\lean{card_zeroTops_le}). The \emph{new candidates} of step~$i$ are the
candidates at~$y_i$ that are not candidates at any $y_{i'}$ with $i'<i$ (\lean{newCands}).

The front reads the bits adaptively: which bits it reads next depends on the bits read so far. The following lemma is the
probabilistic tool for such processes.

\begin{lemma}[Adaptive revealing]\label{lem:revealing}
\formal{Revealing.lean}{RevAgree, IsRevealing, card_le_pow_mul_card_zero, card_le_pow_mul_card_block, card_blocks_le}
Let $I$ be a finite set and $Q\ge1$. For every $i\in\N$ and every bit vector $\xi\in\{0,\dots,Q-1\}^I$ let $R_i(\xi)\subseteq I$
be the set of entries \emph{revealed} in round~$i$. Suppose that the rounds are \emph{adaptive}, $R_i(\xi')=R_i(\xi)$ whenever
$\xi'$ agrees with~$\xi$ on $R_{i'}(\xi)$ for all $i'<i$, and \emph{fresh}, $R_i(\xi)\cap R_{i'}(\xi)=\emptyset$ for $i'<i$. If every
round reveals at most~$\mu$ entries, then for all $a,\ell,j\in\N$ at most
\[
  \bigl(1-Q^{-\ell\mu}\bigr)^j\,Q^{\card I}
\]
bit vectors have the property that each of the $j$ blocks of rounds $[a+b\ell,a+b\ell+\ell)$, $b<j$, contains a round that reveals
an unbroken entry.
\end{lemma}

\begin{outline}
\emph{Zeroing} (\lean{card_le_pow_mul_card_zero}). Let $A$ be a set of bit vectors determined by the rounds before~$i$: if $\xi\in
A$ and $\xi'$ agrees with~$\xi$ on these rounds (\lean{RevAgree}), then $\xi'\in A$. Setting the bits of round~$i$ to~$0$ maps~$A$
into itself and keeps the earlier rounds and round~$i$ itself, by adaptivity and freshness. The vectors with a given image differ
only on round~$i$, so there are at most $Q^\mu$ of them. Hence $\card A\le Q^\mu\card{\{\xi\in A:\text{round $i$ reveals only
broken entries}\}}$, and for $\ell$ consecutive rounds the factor is $Q^{\ell\mu}$ (\lean{card_le_pow_mul_card_block}).
\emph{Blocks} (\lean{card_blocks_le}), by induction on~$j$. Let $E_j$ be the set of bit vectors whose first $j$ blocks each reveal
an unbroken entry; it is determined by the rounds before $a+j\ell$. By zeroing, at least a fraction $Q^{-\ell\mu}$ of~$E_j$ has
block~$j$ all broken, and these vectors are not in $E_{j+1}\subseteq E_j$; so $\card{E_{j+1}}\le(1-Q^{-\ell\mu})\card{E_j}$, and
$\card{E_0}\le Q^{\card I}$.
\end{outline}

\begin{lemma}[The front process]\label{lem:front}
\formal{DisplayTail.lean}{exists_front_eq_top, disp_top, shownSet_subset, card_shownSet_le, isRevealing_newCands, card_newCands_le, card_ucands_succ_lt, card_ucands_add_le, exists_newCands_unbroken, card_running_le}
Suppose that every value has at most $n$ candidates.
\begin{enumerate}[label=(\alph*)]
\item \emph{Closure.} Every element of the shown set is~$0$ or a value~$y_i$ of the front from some $z\in\Sigma\cup\Sigma_0$. Hence,
if no front from a value of $\Sigma\cup\Sigma_0$ runs at step~$N_0$, the shown set has at most $1+\card{\Sigma\cup\Sigma_0}(N_0+1)$
elements.
\item For each~$z$, the new candidates of the steps of the front from~$z$ form an adaptive and fresh revealing scheme with at
most $n$ entries per round. A front that runs at step $a+n$ has an unbroken new candidate at one of the steps $a+1,\dots,a+n$.
\item For all $z$ and~$j$, at most $(1-Q^{-n^2})^jQ^{\card I}$ bit vectors make the front from~$z$ run at step~$jn$.
\end{enumerate}
\end{lemma}

\begin{outline}
(a) The key fact (\lean{exists_front_eq_top}): an unbroken candidate~$e$ at~$y_i$ has its top among the later values of the front.
Indeed $y_{i+1}\le c(e)$, and if $y_{i+1}<c(e)$, then $e$ is still an unbroken candidate at~$y_{i+1}$; induct on $c(e)-y_i$. Now
induct over the rules of the display (\lean{disp_top}). An entry displayed by~(Z) has top~$0$ or a top in~$\Sigma_0$, the
value~$y_0$ of its own front. An entry~$e$ displayed by~(S) through a seed~$z$ has $c(e)=z$ or is an unbroken candidate at
$z=y_0$. An entry~$e$ displayed by~(C) through~$e'$, whose top is $0$ or a front value~$y_i$ by induction, has $c(e)=c(e')$,
or $m(e)=0$ as under~(Z), or is an unbroken candidate at~$y_i$. The seeds are the values~$y_0$ of their own fronts
(\lean{shownSet_subset}). A front that does not run at step~$N_0$ is constant from then on, so it takes at most $N_0+1$ values
(\lean{card_shownSet_le}).
(b) The first $i$ steps of the front read only the bits of the candidates at $y_0,\dots,y_{i-1}$ (\lean{front_congr}), and each of
these is a new candidate at some earlier step (\lean{exists_mem_newCands}); this gives adaptivity, and freshness holds by
definition (\lean{isRevealing_newCands}). The new candidates of a step are candidates at one value, so there are at most~$n$ of
them (\lean{card_newCands_le}). If the front runs at step~$i$ and the new candidates of step $i+1$ are all broken, then the unbroken
candidates at~$y_{i+1}$ form a proper subset of those at~$y_i$ (\lean{card_ucands_succ_lt}): an unbroken candidate~$e$ at~$y_{i+1}$
is not new, so it is a candidate at some $y_{i'}$ with $i'\le i$, and hence at~$y_i$, because $m(e)\le y_{i'}\le y_i\le
y_{i+1}<c(e)$; and the candidate whose top is~$y_{i+1}$ is no longer a candidate at~$y_{i+1}$. A running front has at least one and
at most $n$ unbroken candidates, so $n$ consecutive steps whose new candidates are all broken stop it (\lean{card_ucands_add_le},
\lean{exists_newCands_unbroken}).
(c) If the front runs at step~$jn$, it runs at step $bn+n$ for every $b<j$ (\lean{Running.anti}), and by~(b) each block of steps
$[1+bn,1+bn+n)$ reveals an unbroken entry. \Cref{lem:revealing} with $a=1$ and $\ell=\mu=n$ gives the bound
(\lean{card_running_le}).
\end{outline}

\begin{lemma}[Lemma D: the display is small]\label{lem:D}
\formal{DisplayTail.lean}{card_shownSet_tail}
Suppose that every value has at most $n$ candidates, and let $j\in\N$. At most $(\card\Sigma+n)(1-Q^{-n^2})^jQ^{\card I}$ bit
vectors give a shown set with more than $1+(\card\Sigma+n)(jn+1)$ elements. Equivalently, for uniformly random bits,
\[
  \Pr\bigl[\card{\Ok}>1+(\card\Sigma+n)(jn+1)\bigr]\le(\card\Sigma+n)\bigl(1-Q^{-n^2}\bigr)^j .
\]
\end{lemma}

\begin{outline}
If the shown set is larger, then some front from a value of $\Sigma\cup\Sigma_0$ runs at step~$jn$: otherwise \cref{lem:front}(a)
with $N_0=jn$ bounds the shown set by $1+\card{\Sigma\cup\Sigma_0}(jn+1)$, and $\card{\Sigma\cup\Sigma_0}\le\card\Sigma+n$. A
union bound over the at most $\card\Sigma+n$ starting values, with \cref{lem:front}(c), gives the claim.
\end{outline}

The bound depends neither on the number of entries nor on the depth of the stacks. The seeds are the values~$y_0$ of their own
fronts, so they need no separate term. The composed display has the candidate property and $\card\Sigma\le n$
(\cref{lem:onecandidate}), so every value has at most $n$ candidates (\lean{card_cands_le}) and a time has more than $1+2n(jn+1)$
shown groups with probability at most $2n(1-Q^{-n^2})^j$ (\cref{lem:trunc}); this is what the number of
digits~$P$ of \cref{def:params} is chosen for. The rate $1-Q^{-n^2}$ per block of $n$ steps comes from asking all of the at most
$n^2$ new bits of a block to be broken; it is the reason why $P$, and with it the depth, is so large.

\subsection{The display above a group}\label{sec:displays-view}

Fix a configuration and a value~$d$. An entry~$e$ is \emph{visible} above~$d$ if $d<m(e)\le c(e)$; a candidate at~$d$ has
$m(e)\le d<c(e)$. The \emph{view} above~$d$ consists of the set~$X_d$ of the seeds above~$d$ and the tops of the candidates at~$d$
(\lean{aboveVals}), and the set~$V_d$ of the windows $(m(e),c(e))$ of the unbroken visible entries (\lean{visPairs}). A value is
\emph{reachable} from~$X_d$ through~$V_d$ if it belongs to the least set that contains~$X_d$ and contains~$c$ whenever it contains a
value~$y$ with $m\le y\le c$ for some $(m,c)\in V_d$ (\lean{Reach}).

\begin{lemma}[The display above $d$]\label{lem:view-static}
\formal{DisplayView.lean}{aboveVals, visPairs, Reach, disp_of_cand, mem_dispVals_iff_reach, shown_above_eq}
\begin{enumerate}[label=(\alph*)]
\item A candidate at~$d$ is displayed if its bottom is~$0$, or if it is unbroken and $d$ is a seed.
\item If every candidate at~$d$ is displayed, then a value $y>d$ is displayed if and only if it is reachable from~$X_d$
through~$V_d$.
\item Consider two configurations on the same set of entries, with the same view above~$d$, and in each of which every candidate
at~$d$ is displayed. Then their shown sets have the same elements above~$d$.
\end{enumerate}
\end{lemma}

\begin{outline}
(a) Rules (Z) and~(S), the latter with the seed~$d$ (\lean{disp_of_cand}). (b) (\lean{mem_dispVals_iff_reach}) A seed above~$d$
lies in~$X_d$. The top of a displayed entry~$e$ with $c(e)>d$ is reachable, by induction over the rules: if $m(e)\le d$, then $e$
is a candidate and $c(e)\in X_d$; otherwise $e$ is unbroken and visible, displayed by~(S) through a seed~$z$ or by~(C) through an
entry~$e'$, and $z$, resp.\ $c(e')$, lies in the window of~$e$ and above~$d$, hence is reachable, so that $c(e)$ is reachable in one
more step. Conversely, by induction over reachability: the elements of~$X_d$ are seeds or tops of candidates, and these are
displayed by hypothesis; and if a displayed value lies in the window of an unbroken visible entry~$e$, then $e$ is displayed by~(S)
or~(C). (c) Apply~(b) to both configurations (\lean{shown_above_eq}).
\end{outline}

So, once the candidates at~$d$ are displayed, the display above~$d$ is a function of the view above~$d$. The view is local to the
heads: it consists of their band groups above~$d$, the tops of their candidates at~$d$ and their visible windows. The next
subsection shows that the part of the view belonging to one head changes only at standard events of that head.

\subsection{View invariance}\label{sec:displays-W}

Let $\mathcal S$ be a frame system on $[0,T)$ with a grid $(h_0,r)$, and let $d\in\N$. Its entries, tops, bottoms and band groups are those
of \cref{def:entries}.

\begin{definition}[Standard events and views]\label{def:view}
\formal{ViewInvariance.lean}{FrameSystem.RiseAt, FrameSystem.LastStintAt, FrameSystem.PopsAtZero, FrameSystem.stackEntries, FrameSystem.headAbove, FrameSystem.headVisible}
Let $1\le t<T$ and $g\ge1$.
\begin{enumerate}[label=(\alph*)]
\item There is a \emph{rise} into~$g$ at~$t$ if $t$ is an own time of $J=\inner(t)$, after the start of~$J$ and not a restart
of~$J$, at which the running maximum of~$J$ crosses the floor of~$g$:
$\band(M_J(t-1))<\theta(g)\le\band(M_J(t))$.
\item There is a \emph{last-stint end} at~$t$ at the group~$g$ if, with $J=\inner(t-1)$, $\theta(g)\le\band(M_J(t-1))$ and the
segment of~$J$ current at $t-1$ is never innermost again from~$t$ on: no own time $t'\ge t$ of~$J$ has $s_J(t')=s_J(t-1)$.
\item Rises and last-stint ends are the \emph{standard events}; $t$ is the \emph{time} and $g$ the \emph{group} of the event.
\item $\mathcal S$ \emph{pops at band zero} if every nonempty frame~$F$ other than the root ends at band~$0$: $\beta(f_2-1)=0$.
\item The \emph{view} of~$\mathcal S$ above~$d$ at any time~$t\in\N$ consists of the set $\mathcal A_d(t)$ of the values above~$d$ among the band group
$g(t)$ and the tops~$c_F$ of the entries~$F$ of the stack with $m_F(t)\le d<c_F$, the \emph{candidates} at~$d$
(\lean{FrameSystem.stackEntries}); and the set $\mathcal V_d(t)$ of the pairs $(F,m_F(t))$ for the entries~$F$ of the stack with
$d<m_F(t)\le c_F$, the \emph{visible} entries.
\end{enumerate}
\end{definition}

The frame system of class~U of every walk pops at band zero (\cref{lem:W2b}); so do the frame systems~$\mathcal S_v$ of the heads
(\lean{popsAtZero_headFrames}).

\begin{proposition}[Proposition W: view invariance]\label{prop:W}
\formal{ViewInvariance.lean}{FrameSystem.view_step, FrameSystem.view_interval}
Suppose that $\mathcal S$ pops at band zero and $r<h_0$, and let $p\le u<T$. If no standard event of a group $g\ge d+1$ has its time in
$(p,u]$, then the view above~$d$ is the same at~$p$ and at~$u$: $\mathcal A_d(p)=\mathcal A_d(u)$ and $\mathcal V_d(p)=\mathcal V_d(u)$.
\end{proposition}

\begin{outline}
The statement for $u=p+1$ is \lean{FrameSystem.view_step}, and \lean{FrameSystem.view_interval} follows by induction on~$u$. Fix
$t$ with $1\le t<T$ and no standard event of a group $g\ge d+1$ at~$t$. A frame containing $t-1$ or~$t$ is \emph{common} if it
contains both, \emph{popped} if it contains only $t-1$ (it ends at~$t$), and \emph{pushed} if it contains only~$t$ (it starts
at~$t$); popped and pushed frames lie strictly inside every common frame (\lean{FrameSystem.within_of_contains_not}). Let $C$ be
the innermost common frame (\lean{FrameSystem.InnermostCommon}, \lean{FrameSystem.exists_innermostCommon}; the root is common),
and let $c=G(\band(M_C(t-1)))$.

\emph{Popped and pushed entries.} If some frame is popped at~$t$, the innermost frame at $t-1$ is popped too, so $g(t-1)=0$ because
$\mathcal S$ pops at band zero (\lean{FrameSystem.headGroup_pred_eq_zero}). No entry lies strictly inside a popped entry on the stack at
$t-1$: it would end at~$t$, and its parent, which lies inside the popped entry, would have no own time from~$t$ on. So popped
entries have bottom~$0$ at $t-1$ (\lean{FrameSystem.entryBot_popped}). If some frame is pushed at~$t$, the innermost frame at~$t$
starts at~$t$, so $g(t)=0$; an entry strictly inside a pushed entry and containing~$t$ starts at~$t$ together with its parent, so its
top is the group of a running maximum at the start of a frame, which is~$0$ (\lean{FrameSystem.headGroup_eq_zero_of_push},
\lean{FrameSystem.entryTop_eq_zero}). So pushed entries have bottom~$0$ at~$t$ (\lean{FrameSystem.entryBot_pushed}). Hence popped
and pushed entries are never visible, and they are candidates at~$d$ exactly when their tops exceed~$d$.

\emph{Children of~$C$.} The parent of a popped entry is~$C$: it contains $t-1$, and it contains~$t$ because it has an own time after
the entry ends (\lean{FrameSystem.parent_eq_of_popped}). The parent of a pushed frame is~$C$ or starts at~$t$
(\lean{FrameSystem.parent_of_pushed}), and in the second case the pushed frame has top~$0$. A popped or pushed child of~$C$ has
top~$c$, since the
running maximum of~$C$ is frozen while $C$ is suspended (\lean{FrameSystem.runMax_frozen}), and it is an entry if and only if the
segment of~$C$ current at $t-1$ \emph{returns}, that is, is innermost again at some time $t'\ge t$ (\lean{FrameSystem.Returns},
\lean{FrameSystem.popped_child}, \lean{FrameSystem.pushed_child}).

\emph{Local values.} The \emph{local values} at $t-1$ are the band group $g(t-1)$ and the tops of the entries popped at~$t$; the
local values at~$t$ are $g(t)$ and the tops of the entries pushed at~$t$ (\lean{FrameSystem.LocalPre},
\lean{FrameSystem.LocalPost}). By the above, a positive value is a local value at $t-1$ if and only if it equals~$c$ and $C$ is
innermost at $t-1$ or its segment returns; and it is a local value at~$t$ if and only if either $C$ is innermost at~$t$ and the value
equals $G(\band(M_C(t)))$, or $C$ is not innermost at~$t$, its segment returns, and the value equals~$c$
(\lean{FrameSystem.localPre_iff}, \lean{FrameSystem.localPost_iff}). The key claim
is that the local values at $t-1$ and at~$t$ agree above~$d$ (\lean{FrameSystem.localPre_iff_localPost}). First, if $C$ is innermost
at $t-1$ and its segment does not return, then $c\le d$: otherwise the last stint of that segment ends at~$t$ at the group $c\ge
d+1$, whose floor lies within its band. If $C$ is innermost at~$t$ and restarts there, its band group at~$t$ is~$0$ and its old
segment does not return (\lean{FrameSystem.not_returns_of_restart}), so neither side has a value above~$d$. If $C$ is innermost
at~$t$ without a restart, its segment returns at~$t$ itself, and its group can only rise, from~$c$ to $c'=G(\band(M_C(t)))$; a rise
to $c'>\max(c,d)$ would be a rise into~$c'$ at~$t$, so $c$ and~$c'$ agree above~$d$. If $C$ is not innermost at~$t$, so that some
frame is pushed, both sides are $\{c\}$ when the segment returns; otherwise only the side at $t-1$ can contain~$c$, which is then at
most~$d$.

\emph{Common entries.} The bottom of a common entry at $t-1$, resp.\ at~$t$, is the largest of the tops of the common entries
stacked inside it and of the local values at $t-1$, resp.\ at~$t$ (\lean{FrameSystem.entryBot_pre_le_iff},
\lean{FrameSystem.entryBot_post_le_iff}). Since the local values agree above~$d$, $\max(d,m_F(t-1))=\max(d,m_F(t))$ for every
common entry~$F$ (\lean{FrameSystem.max_entryBot_eq}): it is a candidate at~$d$ at both times or at neither, and if it is visible,
its bottom does not change.

\emph{The view.} A value above~$d$ in $\mathcal A_d(t-1)$ is a local value at $t-1$ (the band group, or the top of a popped entry,
which is a candidate) or the top of a common candidate (\lean{FrameSystem.above_pre_iff}), and likewise at~$t$
(\lean{FrameSystem.above_post_iff}); so $\mathcal A_d(t-1)=\mathcal A_d(t)$. A visible entry at $t-1$ is not popped, since popped
entries have bottom~$0$; so it is common, and visible at~$t$ with the same bottom. Conversely for~$t$, and so
$\mathcal V_d(t-1)=\mathcal V_d(t)$.
\end{outline}

Pushes and pops of entries cost nothing in this proof. At a push, the top of the pushed entry, which is a candidate, equals the band
group of the suspended segment that it replaces; at a pop, the top of the popped entry is replaced by the equal band group of the
resumed segment, unless the running maximum rises at the resumption, which is a rise. What remains are the rises and last-stint
ends above~$d$. A segment rises into a given group at most once and ends its last stint at most once, and on average over the
offset these events are rare (\cref{lem:budgets}).

\subsection{Exposure}\label{sec:displays-exposure}

Let $\mathcal S_0,\dots,\mathcal S_{n-1}$ be frame systems on $[0,T)$ with the composed display at offset~$r$ (\cref{def:composed}), let bits~$\xi$
decide the unbroken entries, and let $d\in\N$. Write $\mathcal A^{(v)}_d(t)$ and $\mathcal V^{(v)}_d(t)$ for the view of~$\mathcal S_v$
above~$d$ at~$t$, and $\Ok_\xi(t)$ for the shown set of the composed display at~$t$.

\begin{proposition}[Exposure]\label{prop:exposure}
\formal{Exposure.lean}{mem_aboveVals_iff, mem_visPairs_iff, views_eq, exposed_of_ne, card_exposed_le}
Let $p,u\in\N$.
\begin{enumerate}[label=(\alph*)]
\item The view of the composed display is assembled from the views of the heads: at every time~$t$, $X_d$ is the union of the
sets $\mathcal A^{(v)}_d(t)$, and $V_d$ consists of the windows $(m,c^{(v)}_F)$ of the unbroken entries $(v,F)$ with
$(F,m)\in\mathcal V^{(v)}_d(t)$. Hence, if every head has the same view above~$d$ at~$p$ and at~$u$, so does the composed display,
whatever the bits.
\item Suppose that $d$ is a seed at~$p$ and at~$u$, and that every head has the same view above~$d$ at~$p$ and at~$u$. If
$\Ok_\xi(p)\cap(d,\infty)\ne\Ok_\xi(u)\cap(d,\infty)$, then some candidate at~$d$, at~$p$ or at~$u$, is broken.
\item Under the hypotheses of~(b), at most $2nQ^{\iota-1}$ of the $Q^\iota$ bit vectors give
$\Ok_\xi(p)\cap(d,\infty)\ne\Ok_\xi(u)\cap(d,\infty)$: for uniformly random bits, the probability is at most $2n/Q$.
\end{enumerate}
\end{proposition}

\begin{outline}
(a) A candidate at~$d$ and a visible entry have nonempty windows, so their frames are entries of the stacks of their heads, and
their bottoms are the bottoms of \cref{def:entries} (\lean{mem_aboveVals_iff}, \lean{mem_visPairs_iff}); the seeds are the band
groups. The bits do not depend on the time (\lean{views_eq}). (b) (\lean{exposed_of_ne}) If every candidate at~$d$, at~$p$ and
at~$u$, is unbroken, then each of them is displayed, $d$ being a seed (\cref{lem:view-static}(a)); by~(a) the two configurations
have the same view above~$d$, and \cref{lem:view-static}(c) gives equal shown sets above~$d$. (c) By the candidate property
(\cref{lem:onecandidate}) there are at most $n$ candidates at~$d$ at each time (\lean{cands}, \lean{card_cands_le}). The bit
vectors that break a given entry are determined by their other $\iota-1$ coordinates, so there are $Q^{\iota-1}$ of them
(\lean{card_some_broken_le}); a union bound over the at most $2n$ candidates gives the claim (\lean{card_exposed_le}).
\end{outline}

In the accounting (\cref{lem:hopterms}), $p<u$ are the ends of an upper hop of a head~$v$, neither lost nor charged, so that no
standard event above its deciding group~$d$ lies in $(p,u]$ and, by \cref{prop:W}, every head has the same view above~$d$ at~$p$
and at~$u$. If the hop is aligned, $d$ is the band group of~$v$ at both ends, a seed, and the proposition applies. If it is not
aligned, $d=0$: every candidate at~$0$ has bottom~$0$ and is displayed by~(Z), so by \cref{lem:view-static}(c) the shown sets
above~$0$ agree for all bits, and the hop is never exposed.
